\documentclass[12pt]{article}
\usepackage[utf8]{inputenc}
\usepackage[T1]{fontenc}
\usepackage{hyperref}
\title{Turkmen Language Resources for K-12 Curriculum Development}
\author{Compiled by Claude Sonnet 5 (claude-sonnet-5) for Andr\'as Kornai}
\date{}
\begin{document}
\maketitle
\section{Introduction}
Turkmen (ISO 639-3 tuk) is the official language of Turkmenistan, a Turkic language with about 16 million speakers, and a language of instruction in schools. The search found mostly materials for foreign learners; I found no national school curriculum or ministry textbook series.
\section{Primary Educational Resources}
\begin{itemize}
\item Tyson and Clark, \emph{Turkmen Language Manual} (Peace Corps, 1993): \url{https://altaica.ru/LIBRARY/Tyson,%20Clark_Turkmen%20Language%20Manual%201993_OCR.pdf}
\item \emph{Beginning Turkmen}, Five College course guides: \url{https://langmedia.fivecolleges.edu/book/export/html/190292}
\item Udemy course, Learn Turkmen: \url{https://www.udemy.com/course/learn-turkmen-beginner-to-advanced/}
\end{itemize}
\section{Linguistic and Reference Materials}
\begin{itemize}
\item Turkmenisch Kurzgrammatik (Landmann, 2013), distributor page: \url{https://isdistribution.com/?p=1410124}
\item UCLA Center for Near Eastern Studies pages on Turkmen: \url{https://www.international.ucla.edu/ccs/article/123625} and \url{https://international.ucla.edu/institute/article/123710}
\item Harvard Caspiana, Turkmenistan academic resources: \url{https://caspiana.omeka.fas.harvard.edu/exhibits/show/turkmenistan/academic-resources}
\end{itemize}
Clark's \emph{Turkmen Reference Grammar} (1998) and \emph{Basic Turkmen} (Dunwoody, 2007) were mentioned in results, but no direct URLs appeared.
\section{Religious and Community Resources}
None found in this search.
\section{Digital and Audio Resources}
\begin{itemize}
\item Turkmen language resources list: \url{https://universeofmemory.com/turkmen-language-resources/}
\item University of Iowa CLCL page: \url{https://clcl.uiowa.edu/node/241}
\end{itemize}
\section{Curriculum Development Context}
Turkmen is already used for instruction, so the main gap is documentation of official school materials, which this search did not find. The listed items are for learners and researchers, not K-12 teaching.
\end{document}